% Recuperación documental para el artículo V; RESULTADO_RECUPERADO.
% Propietario: /Users/ruben/Documents/New project/output/REVISION_EDITORIAL_HMT_MD_20260905/01_FUENTES/INTEGRAL/tree/New project/output/ENSAMBLADO_CAUSAL_RESTAURADO_HMT_MD_20260905/fuente/manuscrito/sections/md/11d_regimenes_cuadraticos_lector_temporal.tex
% Líneas de origen: 4--38, 97--252.
% Sólo se adaptan las etiquetas de referencia y el nivel de títulos.
% La matriz de recuperación conserva la correspondencia y las dependencias.
\section{Regímenes cuadráticos y lectura temporal}
\label{v:rec:chap:regimenes-cuadraticos-tiempo}
\index{álgebra cuadrática!tres regímenes}
\index{cristal de tiempo!realización finita HMT}
\index{tiempo!lector TRIT}

\subsection{Álgebras cuadráticas del TRIT y signaturas de curvatura}
\label{v:rec:sec:regimenes-cuadraticos-recibidos}

La construcción y la demostración primarias de los tres regímenes residen en
la construcción de las álgebras cuadráticas del TRIT,
cuya relación definitoria se presenta en \eqref{eq:algebras-trit}. Para
\(\tau\in\TRIT\), con imagen balanceada
\(\bar\tau\in\{-1,0,+1\}\), se tiene
\[
 \mathbb A_\tau=\RR[X]/(X^2+\bar\tau),
\]
y, si \(J\) es la clase de \(X\),
\[
\begin{array}{ccl}
 \tau=+1&:&J^2=-I,\quad\text{régimen elíptico},\\
 \tau=0&:&J^2=0,\quad\text{régimen parabólico},\\
 \tau=-1&:&J^2=+I,\quad\text{régimen hiperbólico}.
\end{array}
\]
Para los tres valores de \(\tau\), el cociente
\(\mathbb A_\tau=\mathbb R[X]/(X^2+\bar\tau)\) satisface exactamente las tres
identidades cuadráticas anteriores. En particular, la raíz interna de
\(-1\) se realiza como operador en la hoja elíptica; HMT no necesita negar ni
suprimir la extensión compleja. El álgebra parabólica contiene un nilpotente,
y la hiperbólica se descompone mediante los dos idempotentes
\((I\pm J)/2\).

Estas identidades clasifican álgebras y transportes. No seleccionan por sí
solas una métrica espaciotemporal ni prueban que el universo físico ocupe una
de las tres hojas.

\subsection{Lector temporal del TRIT sobre ({-1,0,+1}): apertura, umbral y retorno con memoria}
\label{v:rec:sec:lector-temporal-trit}

Una vez elegida una parametrización ordenada de las transiciones, el lector
temporal activo es
\[
\begin{aligned}
 +1&\longmapsto
 \text{retorno, memoria retenida y pasado},\\
 0&\longmapsto
 \text{umbral de evaluación y presente},\\
 -1&\longmapsto
 \text{apertura bidireccional y futuro}.
\end{aligned}
\]
El lector está \textsc{demostrado con estructura de partida explícita}: sus
conclusiones se obtienen una vez declarados la firma TRIT, el orden de la
trayectoria y la memoria transportada. No convierte el conjunto
\(\{-1,0,+1\}\) en una ecuación de movimiento. El futuro designa el espacio
de prolongaciones compatibles; la supervivencia puede seleccionar una rama
o una multisección sin borrar la genealogía de las restantes.

La frecuencia
\[
 \omega_0=\frac{2\pi}{108\,t_0}
\]
tipa el reloj del retorno completo mediante la entropía de decisión
\(I_{\mathrm{TRIT}}=\log 3\) y el transductor de Boltzmann.
No demuestra, aislada, una fase subarmónica estable. Sí permite construir el
Hamiltoniano finito de referencia que sigue; la selección de un Hamiltoniano
físico robusto requiere condiciones adicionales.

\subsection{Dinámica unitaria en 108 estados: Hamiltoniano, propagador y período}
\label{v:rec:sec:hamiltoniano-reloj-108}

Se fija una orientación \(a\in\{\mathrm{pre},\mathrm{ret}\}\) de las
secciones construidas en \eqref{v:ant:eq:el-secciones-accion} y se escribe
\(\hbar_{\rm int}=\hbar_a\). La misma sección se mantiene en el
Hamiltoniano, el propagador y la conversión térmica posteriores.

Sea
\[
 \mathcal H_{\rm clk}:=\mathbb C^{108},
 \qquad
 \{|j\rangle:j\in\mathbb Z/108\mathbb Z\}
\]
con su producto hermítico canónico. Fijemos
\(\zeta:=\exp(2\pi\mathrm i/108)\), el desplazamiento unitario
\[
 U_{\rm clk}|j\rangle=|j+1\rangle
\]
y la base de Fourier
\[
 |\widetilde k\rangle
 :=\frac1{\sqrt{108}}\sum_{j=0}^{107}\zeta^{jk}|j\rangle,
 \qquad 0\leq k\leq107.
\]
Con \(t_0\in\mathcal L_T^+\) y
\(\hbar_{\rm int}\in\mathcal L_S^+\), definimos
\[
 \boxed{
 H_{\rm clk}
 :=\hbar_{\rm int}\omega_0
 \sum_{k=0}^{107}k\,
 |\widetilde k\rangle\langle\widetilde k|,
 \qquad
 \omega_0=\frac{2\pi}{108\,t_0}.}
 \tag{CLK1}\label{v:rec:eq:hamiltoniano-reloj-108}
\]
El operador tiene tipo energético
\([H_{\rm clk}]=\mathcal L_S\mathcal L_T^{-1}=\mathcal L_E\).

\begin{theorem}[Realización unitaria exacta del calendario \(108\)]
\label{v:rec:thm:reloj-unitario-108}
El operador \(H_{\rm clk}\) es autoadjunto y su propagador durante un paso
satisface
\[
 U_{\rm step}
 :=\exp\!\left(-\frac{\mathrm i}{\hbar_{\rm int}}
 H_{\rm clk}t_0\right)
 =U_{\rm clk},
 \qquad
 U_{\rm step}^{108}=I.
 \tag{CLK2}\label{v:rec:eq:propagador-reloj-108}
\]
Si
\[
 Z_{\rm clk}|j\rangle=\zeta^j|j\rangle,
\]
entonces, para cualquier estado de base \(|j_0\rangle\),
\[
 \langle Z_{\rm clk}\rangle_{n}
 :=\langle j_0|U_{\rm step}^{-n}Z_{\rm clk}
 U_{\rm step}^n|j_0\rangle
 =\zeta^{j_0+n}.
 \tag{CLK3}\label{v:rec:eq:respuesta-reloj-108}
\]
La sucesión posee periodo primitivo \(108\).
\end{theorem}

\begin{proof}
La expresión \eqref{v:rec:eq:hamiltoniano-reloj-108} es una descomposición
espectral con valores reales, luego \(H_{\rm clk}=H_{\rm clk}^\dagger\).
Además,
\[
 U_{\rm clk}|\widetilde k\rangle
 =\zeta^{-k}|\widetilde k\rangle,
\]
mientras que la exponencial de
\(\hbar_{\rm int}\omega_0 k\) durante \(t_0\) produce
\(\exp(-2\pi\mathrm ik/108)=\zeta^{-k}\). Ambos operadores coinciden sobre
una base ortonormal. Las identidades restantes siguen de
\(U_{\rm step}^n|j_0\rangle=|j_0+n\rangle\) y de la primitividad de
\(\zeta\).
\end{proof}

La formulación variacional correspondiente, para una curva
\(\psi\in C^1([t_a,t_b],\mathcal H_{\rm clk})\), no restringida durante la
variación, es
\[
 \boxed{
 L_{\rm clk}(\psi,\dot\psi)
 =\frac{\mathrm i\hbar_{\rm int}}2
 \bigl(
 \langle\psi|\dot\psi\rangle
 -\langle\dot\psi|\psi\rangle
 \bigr)
 -\langle\psi|H_{\rm clk}|\psi\rangle.}
 \tag{CLK4}\label{v:rec:eq:lagrangiano-reloj-108}
\]
El funcional de acción es
\[
 \boxed{
 S_{\rm clk}[\psi]
 :=\int_{t_a}^{t_b}
 L_{\rm clk}(\psi(t),\dot\psi(t))\,\mathrm dt,}
 \tag{CLK5}\label{v:rec:eq:accion-reloj-108}
\]
con \(\delta\psi(t_a)=\delta\psi(t_b)=0\).
\index{Schrödinger!ecuación del reloj finito}
Variar \(\psi\) y \(\bar\psi\) independientemente produce la ecuación de
Schrödinger del reloj finito
\[
 \boxed{\mathrm i\hbar_{\rm int}\dot\psi=H_{\rm clk}\psi}
 \tag{CLK6}\label{v:rec:eq:schrodinger-reloj-108}
\]
junto con su ecuación adjunta \cite{SchrodingerNobel}. Su propagador muestreado
en \(t=nt_0\) es \(\psi_n=U_{\rm step}^{n}\psi_0\). Como \(H_{\rm clk}\) es
autoadjunto, la norma se conserva; un dato inicial normalizado permanece
normalizado. El resultado es exacto respecto de
\(H_{\rm clk},t_0,\hbar_{\rm int}\) declarados y no selecciona un Hamiltoniano
físico ni introduce una frecuencia experimental.

El \cref{v:rec:thm:reloj-unitario-108} demuestra la dinámica unitaria del calendario
respecto de la completación compleja, \(t_0\) y \(\hbar_{\rm int}\)
declarados. Su realización como sistema periódico finito se desarrolla
en la sección~\ref{sec:hmtf2-cristal-temporal}: allí se
fijan el drive, el observable de orden, la subarmonía primitiva de periodo
\(108T_{\rm d}\) y la estabilidad espectral finita. La realización
macroscópica material constituye un reconocimiento externo distinto y no una
laguna del resultado finito.
